% Generated draft: requires mathematical and source review.
\documentclass[11pt]{article}
\usepackage[margin=25mm]{geometry}
\usepackage{fontspec}
\setmainfont{Latin Modern Roman}
\usepackage{amsmath,amssymb,mathtools,unicode-math}
\setmathfont{Latin Modern Math}
\usepackage{enumitem,needspace,xurl,hyperref}
\setlength{\parindent}{0pt}
\setlength{\parskip}{5pt}
\setlength{\emergencystretch}{3em}
\newcommand{\N}{\mathbb{N}}
\newcommand{\Z}{\mathbb{Z}}
\newcommand{\Q}{\mathbb{Q}}
\newcommand{\R}{\mathbb{R}}
\newcommand{\C}{\mathbb{C}}
\newcommand{\F}{\mathbb{F}}
\newcommand{\A}{\mathbb{A}}
\newcommand{\PP}{\mathbb{P}}
\newcommand{\E}{\mathbb{E}}
\newcommand{\eps}{\varepsilon}
\newcommand{\dd}{\,\mathrm{d}}
\begin{document}
\section*{UnsolvedMath ID: 800728}
\textbf{Source problem code:} catalogue entry
\subsection{Definitions and mathematical statement}
A space is called \emph{paracompact} if every open cover admits a locally finite open refinement. Here a family is locally finite if every point has a neighbourhood meeting only finitely many members. Any additional separation requirement in the source's convention for paracompactness remains unspecified.

An algebraic signature \(\Sigma\) specifies operation symbols and their arities; it is uncountable when its set of operation symbols is uncountable. A topological \(\Sigma\)-algebra has continuous interpretations of its operations. To describe the freeness involved, fix a category \(\mathcal C\) of such algebras with continuous homomorphisms. A free object over a space \(X\), when it exists, is a pair \((F_{\mathcal C}(X),\iota_X)\), where \(\iota_X:X\to F_{\mathcal C}(X)\) is continuous and, for every \(A\in\mathcal C\) and every continuous map \(f:X\to A\), there is a unique continuous homomorphism \(\bar f:F_{\mathcal C}(X)\to A\) satisfying \(\bar f\circ\iota_X=f\).

\textbf{General question (Choban).} Under which conditions is the free universal algebra of an uncountable signature over a metrizable space \(X\) paracompact? In the notation above, this asks for conditions ensuring that \(F_{\mathcal C}(X)\) is paracompact. The supplied record does not specify the admissible arities, algebraic identities, category \(\mathcal C\), or separation conventions. Consequently, this notation supplies a framework for the question, but does not determine one particular free-algebra construction.

\textbf{Particular question.} Let \(\R_d\) denote the field of real numbers equipped with the discrete topology. A topological \(\R_d\)-vector space has continuous addition and jointly continuous scalar multiplication \(\R_d\times E\to E\). For a metrizable space \(X\), let \(L(X)\), when it exists in the intended category, be its free topological \(\R_d\)-vector space: there is a continuous map \(\iota_X:X\to L(X)\) such that every continuous map \(f:X\to E\) into an object of that category extends uniquely to a continuous \(\R\)-linear map \(\bar f:L(X)\to E\). Is \(L(X)\) paracompact for every metrizable space \(X\)? The source does not specify whether the vector spaces in this universal property must be Hausdorff; that convention requires review. The discrete topology on the scalar field is an explicit source hypothesis.
\subsection{Short English statement}
Choban asks for conditions ensuring paracompactness of free topological algebras with uncountably many operation symbols over metrizable spaces. The stated special case asks whether every metrizable \(X\) has paracompact free topological linear space \(L(X)\) over the discrete real field. The general algebraic category and relevant separation conventions are unspecified.
\subsection{Sources}
\begin{itemize}
\item Local UnsolvedMath registry, supplied complete source record: permanent ID \(800728\), source problem code \texttt{catalogue entry}, title \emph{Open Problems in Topological Algebra --- Problem 2 (Choban)}. The registry lists Banakh and Choban as proposers, with proposed year 2012, and attributes this problem specifically to Choban.
\item Banakh and Choban, \emph{Some Open Problems in Topological Algebra} (2012), Problem 2. Supplied source URL: \url{https://arxiv.org/abs/1202.1619}. The record reports access on 2026-10-16 and extraction by \texttt{arxiv\_latex}; these are supplied provenance claims, not an independent source visit or verification.
\item The supplied registry labels the problem \texttt{open}. This is a supplied status label, not an independently verified assessment of the literature.
\end{itemize}
\end{document}
